\section{Modelos fundacionales y la visión de la IA general}

\begin{figure}[H]
\centering
\includegraphics[width=0.85\textwidth]{slides-images/slide-77.jpg}
\caption{La idea poderosa que impulsan los Foundation Models\protect\footnotemark}
\end{figure}
\footnotetext{Diapositivas, página 77.}

Al final de la clase, Amini amplía la mirada hacia la visión más amplia de los Foundation Models:

\begin{importantbox}{La visión central de los Foundation Models}
\begin{itemize}
    \item \textbf{Sistema central de razonamiento}: ¿pueden los modelos fundacionales generativos ofrecer un motor de razonamiento unificado para la IA general?
    \item \textbf{IA que diseña IA}: ¿puede usarse la IA para mejorar y hacer evolucionar a la propia IA?
    \item \textbf{IA generativa multidominio}: de las imágenes a la biología, al lenguaje y más allá; con un enorme potencial, pero que exige cautela
    \item \textbf{La relación entre la inteligencia artificial y la inteligencia humana}: entender los vínculos y las diferencias entre ambas
\end{itemize}
\end{importantbox}

Estas preguntas no son solo desafíos técnicos, sino que también implican una reflexión científica y filosófica profunda. Como enfatiza Amini: el poder y la cautela van de la mano (power and caution); al mismo tiempo que impulsamos la frontera de la IA, debemos considerar con prudencia su impacto social.

%% =====================================================
